% Nguồn: Tuyển tập đề tham khảo môn Toán TS10 2025-2026 – TP. Hồ Chí Minh
% Đợt 5: Quận 5 – Đề 1 và Đề 2 (14 câu tự luận)

% ==============================================================================
% QUẬN 5 – ĐỀ 1
% ==============================================================================

\begin{bt}%[9D6H18-2]
% nguon: Đề tham khảo Tuyển sinh 10 TP.HCM 2025-2026 – Quận 5 – Đề 1 – Bài 1
Cho parabol $(P)\colon y = -\dfrac{1}{2}x^2$.
\begin{enumerate}[a)]
    \item Vẽ $(P)$ trên hệ trục tọa độ $Oxy$.
    \item Tìm tọa độ các điểm $A$ thuộc $(P)$ có hoành độ gấp hai lần tung độ.
\end{enumerate}
\loigiai{
\begin{enumerate}[a)]
    \item Bảng giá trị tương ứng của $x$ và $y$:
    \begin{center}
    \begin{tabular}{|c|c|c|c|c|c|}
        \hline
        $x$ & $-4$ & $-2$ & $0$ & $2$ & $4$ \\ \hline
        $y = -\dfrac{1}{2}x^2$ & $-8$ & $-2$ & $0$ & $-2$ & $-8$ \\ \hline
    \end{tabular}
    \end{center}
    Đồ thị $(P)$ là một parabol đi qua gốc tọa độ $O(0;0)$, nhận trục $Oy$ làm trục đối xứng và nằm phía dưới trục hoành.
    \item Gọi điểm $A(x_A; y_A) \in (P)$ thỏa mãn hoành độ gấp hai lần tung độ, tức là $x_A = 2y_A \Leftrightarrow y_A = \dfrac{x_A}{2}$.\\
    Vì $A \in (P)$ nên ta có phương trình:
    $\dfrac{x_A}{2} = -\dfrac{1}{2}x_A^2 \Leftrightarrow x_A^2 + x_A = 0 \Leftrightarrow x_A(x_A + 1) = 0 \Leftrightarrow \left[\begin{aligned} &x_A = 0 \\ &x_A = -1. \end{aligned}\right.$\\
    Với $x_A = 0 \Rightarrow y_A = 0$, ta được điểm $A_1(0; 0)$.\\
    Với $x_A = -1 \Rightarrow y_A = -\dfrac{1}{2}$, ta được điểm $A_2\left(-1; -\dfrac{1}{2}\right)$.\\
    Vậy tọa độ các điểm cần tìm là $(0; 0)$ và $\left(-1; -\dfrac{1}{2}\right)$.
\end{enumerate}
}
\end{bt}

\begin{bt}%[9D6H20-2]
% nguon: Đề tham khảo Tuyển sinh 10 TP.HCM 2025-2026 – Quận 5 – Đề 1 – Bài 2
Cho phương trình $x^2 + 2x - 2025 = 0$.
\begin{enumerate}[a)]
    \item Chứng minh phương trình có hai nghiệm phân biệt trái dấu.
    \item Gọi $x_1$, $x_2$ là hai nghiệm của phương trình. Không giải phương trình, tính giá trị của biểu thức $A = \dfrac{(1 - x_1)x_2}{x_1} + \dfrac{(1 - x_2)x_1}{x_2}$.
\end{enumerate}
\loigiai{
\begin{enumerate}[a)]
    \item Phương trình $x^2 + 2x - 2025 = 0$ có hệ số $a = 1$ và $c = -2025$.\\
    Ta có tích $a \cdot c = 1 \cdot (-2025) = -2025 < 0$.\\
    Do $a$ và $c$ trái dấu nên phương trình luôn có hai nghiệm phân biệt trái dấu $x_1, x_2$.
    \item Theo định lí Viète, ta có:
    $\left\{\begin{aligned} &S = x_1 + x_2 = -2 \\ &P = x_1 x_2 = -2025. \end{aligned}\right.$\\
    Biến đổi biểu thức $A$:
    $\begin{aligned}
    A &= \dfrac{(1 - x_1)x_2}{x_1} + \dfrac{(1 - x_2)x_1}{x_2} = \dfrac{x_2^2(1 - x_1) + x_1^2(1 - x_2)}{x_1 x_2} \\
    &= \dfrac{(x_1^2 + x_2^2) - x_1 x_2(x_1 + x_2)}{x_1 x_2} = \dfrac{(x_1 + x_2)^2 - 2x_1 x_2 - x_1 x_2(x_1 + x_2)}{x_1 x_2} \\
    &= \dfrac{S^2 - 2P - P \cdot S}{P} = \dfrac{(-2)^2 - 2(-2025) - (-2025)(-2)}{-2025} \\
    &= \dfrac{4 + 4050 - 4050}{-2025} = -\dfrac{4}{2025}.
    \end{aligned}$\\
    Vậy $A = -\dfrac{4}{2025}$.
\end{enumerate}
}
\end{bt}

\begin{bt}%[9D8H26-3]
% nguon: Đề tham khảo Tuyển sinh 10 TP.HCM 2025-2026 – Quận 5 – Đề 1 – Bài 3
Một bó hoa gồm 5 bông hoa màu đỏ và 3 bông hoa màu vàng. Bạn Linh chọn ngẫu nhiên cùng lúc 4 bông hoa từ bó hoa đó.
\begin{enumerate}[a)]
    \item Liệt kê các cách chọn mà bạn Linh có thể thực hiện.
    \item Tính xác suất của mỗi biến cố sau:
    \begin{itemize}
        \item $A$: “Trong 4 bông hoa được chọn ra, có đúng 1 bông hoa màu đỏ”.
        \item $B$: “Trong 4 bông hoa được chọn ra, có ít nhất 1 bông hoa màu đỏ”.
    \end{itemize}
\end{enumerate}
\loigiai{
\begin{enumerate}[a)]
    \item Tổng số bông hoa trong bó là $5 + 3 = 8$ bông hoa.\\
    Khi chọn ngẫu nhiên cùng lúc 4 bông hoa, các khả năng về số bông hoa theo màu có thể xảy ra là:
    \begin{itemize}
        \item Chọn 1 bông màu đỏ và 3 bông màu vàng.
        \item Chọn 2 bông màu đỏ và 2 bông màu vàng.
        \item Chọn 3 bông màu đỏ và 1 bông màu vàng.
        \item Chọn 4 bông màu đỏ và 0 bông màu vàng.
    \end{itemize}
    Tổng số kết quả có thể xảy ra khi chọn 4 bông hoa từ 8 bông hoa là $C_8^4 = \dfrac{8 \cdot 7 \cdot 6 \cdot 5}{4 \cdot 3 \cdot 2 \cdot 1} = 70$ cách chọn.
    \item Tính xác suất:
    \begin{itemize}
        \item Biến cố $A$: “Trong 4 bông hoa được chọn ra, có đúng 1 bông hoa màu đỏ”.\\
        Để chọn được đúng 1 bông hoa đỏ, ta chọn 1 bông đỏ từ 5 bông đỏ (có 5 cách) và 3 bông vàng từ 3 bông vàng (có 1 cách).\\
        Số kết quả thuận lợi cho biến cố $A$ là $n(A) = 5 \cdot 1 = 5$.\\
        Xác suất của biến cố $A$ là $P(A) = \dfrac{n(A)}{70} = \dfrac{5}{70} = \dfrac{1}{14}$.
        \item Biến cố $B$: “Trong 4 bông hoa được chọn ra, có ít nhất 1 bông hoa màu đỏ”.\\
        Biến cố đối của $B$ là $\overline{B}$: “Trong 4 bông hoa được chọn ra, không có bông hoa màu đỏ nào” (tức cả 4 bông đều màu vàng).\\
        Vì bó hoa chỉ có 3 bông màu vàng nên không thể chọn được 4 bông màu vàng, tức là $n(\overline{B}) = 0$.\\
        Do đó biến cố $B$ là biến cố chắc chắn, nên $P(B) = 1$.
    \end{itemize}
\end{enumerate}
}
\end{bt}

\begin{bt}%[9D6V21-3]
% nguon: Đề tham khảo Tuyển sinh 10 TP.HCM 2025-2026 – Quận 5 – Đề 1 – Bài 4
Nhà ông Hiền có một mảnh đất hình chữ nhật có chiều rộng là $5\text{ m}$ và chiều dài là $50\text{ m}$. Do việc mở rộng đường, nên nhà nước đã phải thu hồi một phần diện tích đất của nhà ông Hiền (phần hình tam giác vuông ở góc mảnh đất, với $0 < x < 5$).
\begin{enumerate}[a)]
    \item Viết và thu gọn biểu thức $A$ biểu thị theo $x$ diện tích phần đất bị thu hồi của nhà ông Hiền, biết hai cạnh góc vuông của phần tam giác bị thu hồi lần lượt là $5 - x$ (m) và $50 - 2x$ (m).
    \item Phần đất bị thu hồi của nhà ông Hiền được nhà nước đền bù số tiền $440$ triệu đồng. Tìm giá trị $x$, biết giá đền bù đất bị thu hồi là $10\text{ triệu đồng/m}^2$.
\end{enumerate}
\loigiai{
\begin{enumerate}[a)]
    \item Phần đất bị thu hồi là một hình tam giác vuông có độ dài hai cạnh góc vuông lần lượt là $5 - x\text{ (m)}$ và $50 - 2x\text{ (m)}$.\\
    Diện tích phần đất bị thu hồi được biểu diễn bởi:
    $A = \dfrac{1}{2}(5 - x)(50 - 2x) = (5 - x)(25 - x) = x^2 - 30x + 125\text{ (m}^2\text{)}$.
    \item Diện tích phần đất bị thu hồi tương ứng với số tiền đền bù là:
    $S = \dfrac{440}{10} = 44\text{ (m}^2\text{)}$.\\
    Theo câu a, ta có phương trình:
    $x^2 - 30x + 125 = 44 \Leftrightarrow x^2 - 30x + 81 = 0 \Leftrightarrow (x - 3)(x - 27) = 0 \Leftrightarrow \left[\begin{aligned} &x = 3\text{ (thỏa mãn } 0 < x < 5) \\ &x = 27\text{ (loại vì } x > 5). \end{aligned}\right.$\\
    Vậy giá trị cần tìm là $x = 3$.
\end{enumerate}
}
\end{bt}

\begin{bt}%[9H0V31-3]
% nguon: Đề tham khảo Tuyển sinh 10 TP.HCM 2025-2026 – Quận 5 – Đề 1 – Bài 5
Một hộp sữa đặc dạng hình trụ có chiều cao bằng $12\text{ cm}$. Biết thể tích của hộp sữa là $192\pi\text{ cm}^3$.
\begin{enumerate}[a)]
    \item Tính diện tích của vỏ hộp sữa (kể cả hai nắp hộp; bỏ qua các nếp gấp cũng như các phần gắn thêm).
    \item Tính số vỏ hộp sữa nhiều nhất có thể làm được từ miếng kim loại có diện tích là $10\text{ m}^2$ (hao phí không đáng kể).
\end{enumerate}
\loigiai{
\begin{enumerate}[a)]
    \item Gọi $r$ là bán kính đáy của hộp sữa đặc ($r > 0$, đơn vị cm).\\
    Thể tích hình trụ được tính theo công thức $V = \pi r^2 h$.\\
    Theo đề bài:
    $192\pi = \pi \cdot r^2 \cdot 12 \Leftrightarrow r^2 = \dfrac{192}{12} = 16 \Rightarrow r = 4\text{ cm}$.\\
    Diện tích toàn phần của vỏ hộp sữa (diện tích vỏ hộp kể cả hai nắp) là:
    $S_{\text{tp}} = 2\pi r h + 2\pi r^2 = 2\pi \cdot 4 \cdot 12 + 2\pi \cdot 4^2 = 96\pi + 32\pi = 128\pi\text{ cm}^2 \approx 402{,}12\text{ cm}^2$.
    \item Đổi diện tích miếng kim loại ra $\text{cm}^2$:
    $10\text{ m}^2 = 10 \cdot 10\,000 = 100\,000\text{ cm}^2$.\\
    Số vỏ hộp sữa nhiều nhất có thể sản xuất được là:
    $N = \left\lfloor \dfrac{100\,000}{128\pi} \right\rfloor \approx \lfloor 248{,}68 \rfloor = 248\text{ (vỏ hộp)}$.\\
    Vậy có thể làm được nhiều nhất 248 vỏ hộp sữa.
\end{enumerate}
}
\end{bt}

\begin{bt}%[9D1H3-4]
% nguon: Đề tham khảo Tuyển sinh 10 TP.HCM 2025-2026 – Quận 5 – Đề 1 – Bài 6
Nhân dịp tết Nguyên đán, một trường trung học cơ sở đã tổ chức cho học sinh các khối lớp 6, 7, 8 và 9 quyên góp tiền ủng hộ tự nguyện nhằm góp phần nhỏ giúp đỡ những học sinh có hoàn cảnh khó khăn. Kết thúc đợt quyên góp, bình quân mỗi học sinh của khối lớp 9 quyên góp được $22\,000$ đồng; bình quân mỗi học sinh của ba khối lớp 6, 7 và 8 quyên được $18\,000$ đồng. Hỏi bình quân mỗi học sinh của trường quyên góp được bao nhiêu tiền? Biết tổng số học sinh của ba khối lớp 6, 7 và 8 gấp 3 lần tổng số của học sinh khối lớp 9.
\loigiai{
Gọi $x$ là tổng số học sinh khối lớp 9 ($x \in \mathbb{N}^*$).\\
Vì tổng số học sinh của ba khối 6, 7 và 8 gấp 3 lần số học sinh khối 9 nên số học sinh ba khối 6, 7 và 8 là $3x$ (học sinh).\\
Tổng số học sinh toàn trường là: $x + 3x = 4x$ (học sinh).\\
Tổng số tiền học sinh khối lớp 9 quyên góp được là:
$22\,000 \cdot x$ (đồng).\\
Tổng số tiền học sinh ba khối 6, 7 và 8 quyên góp được là:
$18\,000 \cdot 3x = 54\,000x$ (đồng).\\
Tổng số tiền toàn trường quyên góp được là:
$22\,000x + 54\,000x = 76\,000x$ (đồng).\\
Bình quân mỗi học sinh toàn trường quyên góp được số tiền là:
$\dfrac{76\,000x}{4x} = 19\,000$ (đồng).\\
Vậy bình quân mỗi học sinh của trường quyên góp được $19\,000$ đồng.
}
\end{bt}

\begin{bt}%[9H9C29-3]
% nguon: Đề tham khảo Tuyển sinh 10 TP.HCM 2025-2026 – Quận 5 – Đề 1 – Bài 7
Cho tam giác $ABC$ vuông tại $A$ ($AB > AC$), đường cao $AH$; gọi $O$ là tâm đường tròn ngoại tiếp của tam giác. Các tiếp tuyến với đường tròn $(O)$ tại $A$ và tại $B$ cắt nhau tại $M$. Nối $CM$ cắt $AH$ tại $I$. Nối $OM$ cắt $AB$ tại $J$.
\begin{enumerate}[a)]
    \item Chứng minh các tam giác $MOB$ và $CAH$ đồng dạng.
    \item Chứng minh $I$ là trung điểm của $AH$.
    \item Cho $BC = 2R$ và $OM = x$ ($R > 0; x > 0$). Tính $AB$, $AH$ theo $R$ và $x$.
\end{enumerate}
\loigiai{
\begin{enumerate}[a)]
    \item Vì tam giác $ABC$ vuông tại $A$ nên tâm $O$ đường tròn ngoại tiếp tam giác là trung điểm của cạnh huyền $BC$.\\
    Ta có $MA, MB$ là hai tiếp tuyến cắt nhau tại $M$ của đường tròn $(O)$, suy ra $MA = MB$ và $OM$ là tia phân giác của $\widehat{AOB}$.\\
    Do đó $\widehat{BOM} = \dfrac{1}{2}\widehat{AOB}$.\\
    Mặt khác, góc nội tiếp $\widehat{ACB}$ chắn cung nhỏ $AB$ nên $\widehat{ACB} = \dfrac{1}{2}\widehat{AOB}$.\\
    Suy ra $\widehat{BOM} = \widehat{ACB} = \widehat{ACH}$.\\
    Xét $\triangle MOB$ và $\triangle CAH$ có:
    $\left\{\begin{aligned} &\widehat{OBM} = \widehat{AHC} = 90^\circ \\ &\widehat{MOB} = \widehat{ACH}\text{ (chứng minh trên)}. \end{aligned}\right.$\\
    Vậy $\triangle MOB \sim \triangle CAH$ (g-g).
    \item Gọi $K$ là giao điểm của đường thẳng $CA$ và tiếp tuyến $MB$.\\
    Vì $MB$ là tiếp tuyến tại $B$ nên $MB \perp BC$. Mặt khác $AH \perp BC$ nên $AH \parallel BK$.\\
    Tam giác $KBC$ vuông tại $B$ có $AB \perp KC$ (vì $\widehat{BAC} = 90^\circ$).\\
    Ta có $MA = MB$ nên tam giác $MAB$ cân tại $M$, do đó $\widehat{MAB} = \widehat{MBA}$.\\
    Lại có $\widehat{MAB} + \widehat{MAK} = \widehat{CAB} = 90^\circ$ và $\widehat{MBA} + \widehat{MKA} = 90^\circ$ (trong tam giác vuông $KAB$).\\
    Suy ra $\widehat{MAK} = \widehat{MKA}$, do đó tam giác $MAK$ cân tại $M \Rightarrow MA = MK$.\\
    Từ đó suy ra $MB = MK$, tức là $M$ là trung điểm của đoạn thẳng $BK$.\\
    Trong tam giác $CKB$, do $AH \parallel BK$, áp dụng hệ quả định lí Thales ta có:
    $\dfrac{IH}{MB} = \dfrac{CI}{CM} = \dfrac{IA}{MK}$.\\
    Vì $MB = MK$ nên $IH = IA$, suy ra $I$ là trung điểm của $AH$.
    \item Xét tam giác $OBM$ vuông tại $B$ có bán kính $OB = R$, $OM = x$ ($x > R$):\\
    Ta có $MA = MB$ và $OA = OB = R$ nên $OM$ là đường trung trực của đoạn $AB$, suy ra $OM \perp AB$ tại trung điểm $J$ của $AB$.\\
    Áp dụng hệ thức lượng trong tam giác vuông $OBM$ có đường cao $BJ$:
    $\begin{aligned}
    &BM = \sqrt{OM^2 - OB^2} = \sqrt{x^2 - R^2}, \\
    &BJ \cdot OM = OB \cdot BM \Rightarrow BJ = \dfrac{OB \cdot BM}{OM} = \dfrac{R\sqrt{x^2 - R^2}}{x}.
    \end{aligned}$\\
    Vì $J$ là trung điểm của $AB$ nên $AB = 2BJ = \dfrac{2R\sqrt{x^2 - R^2}}{x}$.\\
    Tam giác $ABC$ vuông tại $A$ có cạnh huyền $BC = 2R$, áp dụng định lí Pythagore:
    $AC = \sqrt{BC^2 - AB^2} = \sqrt{4R^2 - \dfrac{4R^2(x^2 - R^2)}{x^2}} = \sqrt{\dfrac{4R^4}{x^2}} = \dfrac{2R^2}{x}$.\\
    Đường cao $AH$ của tam giác vuông $ABC$ được tính bởi:
    $AH = \dfrac{AB \cdot AC}{BC} = \dfrac{\dfrac{2R\sqrt{x^2 - R^2}}{x} \cdot \dfrac{2R^2}{x}}{2R} = \dfrac{2R^2\sqrt{x^2 - R^2}}{x^2}$.\\
    Vậy $AB = \dfrac{2R\sqrt{x^2 - R^2}}{x}$ và $AH = \dfrac{2R^2\sqrt{x^2 - R^2}}{x^2}$.
\end{enumerate}
}
\end{bt}

% ==============================================================================
% QUẬN 5 – ĐỀ 2
% ==============================================================================

\begin{bt}%[9D6H18-2]
% nguon: Đề tham khảo Tuyển sinh 10 TP.HCM 2025-2026 – Quận 5 – Đề 2 – Bài 1
Cho parabol $(P)\colon y = -\dfrac{1}{2}x^2$.
\begin{enumerate}[a)]
    \item Vẽ đồ thị $(P)$.
    \item Tìm các điểm $M$ thuộc $(P)$ sao cho tung độ bằng $-8$.
\end{enumerate}
\loigiai{
\begin{enumerate}[a)]
    \item Bảng giá trị tương ứng của $x$ và $y$:
    \begin{center}
    \begin{tabular}{|c|c|c|c|c|c|}
        \hline
        $x$ & $-4$ & $-2$ & $0$ & $2$ & $4$ \\ \hline
        $y = -\dfrac{1}{2}x^2$ & $-8$ & $-2$ & $0$ & $-2$ & $-8$ \\ \hline
    \end{tabular}
    \end{center}
    Đồ thị $(P)$ là một parabol đi qua gốc tọa độ $O(0;0)$, nhận trục $Oy$ làm trục đối xứng và bề lõm quay xuống dưới.
    \item Điểm $M(x_M; y_M) \in (P)$ có tung độ $y_M = -8$.\\
    Thay $y = -8$ vào phương trình parabol $(P)$, ta được:
    $-8 = -\dfrac{1}{2}x^2 \Leftrightarrow x^2 = 16 \Leftrightarrow x = \pm 4$.\\
    Vậy có hai điểm thỏa mãn yêu cầu bài toán là $M_1(4; -8)$ và $M_2(-4; -8)$.
\end{enumerate}
}
\end{bt}

\begin{bt}%[9D6H20-2]
% nguon: Đề tham khảo Tuyển sinh 10 TP.HCM 2025-2026 – Quận 5 – Đề 2 – Bài 2
Cho phương trình $x^2 - 4x - 6 = 0$.
\begin{enumerate}[a)]
    \item Chứng minh phương trình trên có hai nghiệm phân biệt $x_1$, $x_2$.
    \item Không giải phương trình, hãy tính giá trị của biểu thức: $A = \dfrac{x_1 x_2}{4 - x_1} + \dfrac{x_1 x_2}{4 - x_2}$.
\end{enumerate}
\loigiai{
\begin{enumerate}[a)]
    \item Phương trình $x^2 - 4x - 6 = 0$ có các hệ số $a = 1$ và $c = -6$.\\
    Vì $a \cdot c = 1 \cdot (-6) = -6 < 0$ nên phương trình luôn có hai nghiệm phân biệt $x_1, x_2$ trái dấu nhau.
    \item Theo định lí Viète, ta có:
    $\left\{\begin{aligned} &x_1 + x_2 = 4 \\ &x_1 x_2 = -6. \end{aligned}\right.$\\
    Từ $x_1 + x_2 = 4$, suy ra $4 - x_1 = x_2$ và $4 - x_2 = x_1$.\\
    Do đó biểu thức $A$ có thể viết lại thành:
    $A = \dfrac{x_1 x_2}{x_2} + \dfrac{x_1 x_2}{x_1} = x_1 + x_2 = 4$.\\
    Vậy $A = 4$.
\end{enumerate}
}
\end{bt}

\begin{bt}%[9D8H26-3]
% nguon: Đề tham khảo Tuyển sinh 10 TP.HCM 2025-2026 – Quận 5 – Đề 2 – Bài 3
Trong một bó hoa có tất cả 5 bông hoa màu đỏ và 3 bông hoa màu vàng. Bạn Linh chọn ngẫu nhiên đồng thời 4 bông hoa từ bó hoa ấy.
\begin{enumerate}[a)]
    \item Liệt kê các cách chọn mà bạn Linh có thể thực hiện.
    \item Tính xác suất của mỗi biến cố sau:
    \begin{itemize}
        \item $A$: “Trong 4 bông hoa được chọn ra, có đúng 1 bông hoa màu đỏ”.
        \item $B$: “Trong 4 bông hoa được chọn ra, có ít nhất 1 bông hoa màu đỏ”.
    \end{itemize}
\end{enumerate}
\loigiai{
\begin{enumerate}[a)]
    \item Tổng số bông hoa là $5 + 3 = 8$ bông hoa.\\
    Khi chọn cùng lúc 4 bông hoa, các trường hợp về số bông theo từng màu gồm có:
    \begin{itemize}
        \item Trường hợp 1: Chọn 1 bông màu đỏ và 3 bông màu vàng.
        \item Trường hợp 2: Chọn 2 bông màu đỏ và 2 bông màu vàng.
        \item Trường hợp 3: Chọn 3 bông màu đỏ và 1 bông màu vàng.
        \item Trường hợp 4: Chọn 4 bông màu đỏ và 0 bông màu vàng.
    \end{itemize}
    Số phần tử của không gian mẫu là $n(\Omega) = C_8^4 = 70$.
    \item Tính xác suất:
    \begin{itemize}
        \item Biến cố $A$: “Trong 4 bông hoa được chọn ra, có đúng 1 bông hoa màu đỏ”.\\
        Số cách chọn 1 bông đỏ từ 5 bông đỏ và 3 bông vàng từ 3 bông vàng là $n(A) = C_5^1 \cdot C_3^3 = 5 \cdot 1 = 5$.\\
        Xác suất của biến cố $A$ là $P(A) = \dfrac{5}{70} = \dfrac{1}{14}$.
        \item Biến cố $B$: “Trong 4 bông hoa được chọn ra, có ít nhất 1 bông hoa màu đỏ”.\\
        Xét biến cố đối $\overline{B}$: “Chọn được 4 bông hoa không có bông đỏ nào” (tức cả 4 bông đều màu vàng).\\
        Do chỉ có 3 bông màu vàng nên số cách chọn 4 bông hoa màu vàng là $0$, tức $n(\overline{B}) = 0$.\\
        Suy ra biến cố $B$ chắc chắn xảy ra, vậy xác suất $P(B) = 1$.
    \end{itemize}
\end{enumerate}
}
\end{bt}

\begin{bt}%[9D1V3-4]
% nguon: Đề tham khảo Tuyển sinh 10 TP.HCM 2025-2026 – Quận 5 – Đề 2 – Bài 4
Một xí nghiệp dự tính vận chuyển hàng bằng hai xe tải và đang phân vân giữa hai phương án là mua hẳn hai chiếc xe tải hoặc thuê hai xe tải (hai xe giống nhau). Nếu mua hai xe với giá mỗi xe là $200\,000\,000$ đồng thì mỗi ngày xí nghiệp phải trả tiền công cho tài xế và mua nhiên liệu cho một xe là $5\,000\,000$ đồng. Còn nếu thuê xe thì giá thuê một xe tải là $10\,000\,000$ đồng/ngày (đã bao gồm tiền công cho tài xế và nhiên liệu).
\begin{enumerate}[a)]
    \item Gọi $T$ là tổng số tiền xí nghiệp bỏ ra để vận chuyển hàng sau $n$ ngày. Lập công thức tính $T$ theo $n$ đối với mỗi phương án.
    \item Hỏi xí nghiệp vận chuyển hàng từ bao nhiêu ngày trở lên thì nên chọn phương án mua xe sẽ tiết kiệm hơn phương án thuê xe?
\end{enumerate}
\loigiai{
\begin{enumerate}[a)]
    \item Lập công thức tính tổng chi phí $T$ sau $n$ ngày vận chuyển:
    \begin{itemize}
        \item Phương án mua 2 xe tải:\\
        Chi phí mua 2 xe tải ban đầu là:
        $2 \cdot 200\,000\,000 = 400\,000\,000$ (đồng).\\
        Chi phí tài xế và nhiên liệu cho cả 2 xe mỗi ngày là:
        $2 \cdot 5\,000\,000 = 10\,000\,000$ (đồng/ngày).\\
        Tổng số tiền chi sau $n$ ngày theo phương án mua xe là:
        $T_1(n) = 400\,000\,000 + 10\,000\,000n$ (đồng).
        \item Phương án thuê 2 xe tải:\\
        Giá thuê cả 2 xe mỗi ngày là:
        $2 \cdot 10\,000\,000 = 20\,000\,000$ (đồng/ngày).\\
        Tổng số tiền chi sau $n$ ngày theo phương án thuê xe là:
        $T_2(n) = 20\,000\,000n$ (đồng).
    \end{itemize}
    \item Để phương án mua xe tiết kiệm hơn phương án thuê xe thì tổng chi phí mua xe phải nhỏ hơn tổng chi phí thuê xe:
    $T_1(n) < T_2(n) \Leftrightarrow 400\,000\,000 + 10\,000\,000n < 20\,000\,000n \Leftrightarrow 10\,000\,000n > 400\,000\,000 \Leftrightarrow n > 40$.\\
    Vì $n$ là số ngày nguyên dương nên $n \ge 41$.\\
    Vậy xí nghiệp vận chuyển hàng từ 41 ngày trở lên thì nên chọn phương án mua xe sẽ tiết kiệm hơn phương án thuê xe.
\end{enumerate}
}
\end{bt}

\begin{bt}%[9H0V31-3]
% nguon: Đề tham khảo Tuyển sinh 10 TP.HCM 2025-2026 – Quận 5 – Đề 2 – Bài 5
Một xe tải chở hàng có thùng xe dạng hình hộp chữ nhật với kích thước: chiều dài $3\text{ m}$, chiều rộng $2\text{ m}$ và chiều cao $1{,}5\text{ m}$.
\begin{enumerate}[a)]
    \item Tính thể tích của thùng xe (độ dày của thùng xe xem như không đáng kể).
    \item Người ta xếp vào thùng xe tải các thùng hàng loại A có dạng hình lập phương độ dài cạnh $40\text{ cm}$ để vận chuyển giao hàng. Hỏi mỗi lần vận chuyển xe có thể chở được nhiều nhất bao nhiêu thùng hàng loại A?
\end{enumerate}
\loigiai{
\begin{enumerate}[a)]
    \item Thể tích của thùng xe tải là:
    $V = 3 \cdot 2 \cdot 1{,}5 = 9\text{ (m}^3\text{)}$.
    \item Đổi độ dài cạnh của mỗi thùng hàng loại A: $40\text{ cm} = 0{,}4\text{ m}$.\\
    Vì các thùng hàng là hình lập phương và phải được xếp vừa vặn vào trong thùng xe theo từng hàng dọc, ngang và cao nên:
    \begin{itemize}
        \item Số thùng hàng xếp được nhiều nhất theo chiều dài là:
        $\left\lfloor \dfrac{3}{0{,}4} \right\rfloor = \lfloor 7{,}5 \rfloor = 7\text{ (thùng)}$.
        \item Số thùng hàng xếp được nhiều nhất theo chiều rộng là:
        $\left\lfloor \dfrac{2}{0{,}4} \right\rfloor = \lfloor 5 \rfloor = 5\text{ (thùng)}$.
        \item Số thùng hàng xếp được nhiều nhất theo chiều cao là:
        $\left\lfloor \dfrac{1{,}5}{0{,}4} \right\rfloor = \lfloor 3{,}75 \rfloor = 3\text{ (thùng)}$.
    \end{itemize}
    Vậy mỗi lần vận chuyển xe tải có thể chở được nhiều nhất số thùng hàng loại A là:
    $7 \cdot 5 \cdot 3 = 105\text{ (thùng)}$.
\end{enumerate}
}
\end{bt}

\begin{bt}%[9D1V3-4]
% nguon: Đề tham khảo Tuyển sinh 10 TP.HCM 2025-2026 – Quận 5 – Đề 2 – Bài 6
Ông Năm dự định cải tạo một mảnh vườn hình chữ nhật có chiều dài bằng $2{,}5$ lần chiều rộng. Ông nhận thấy rằng, nếu đào một cái hồ có mặt hồ là hình chữ nhật thì sẽ chiếm mất $3\%$ diện tích mảnh vườn. Còn nếu giảm chiều dài mặt hồ $5\text{ m}$ và tăng chiều rộng mặt hồ $2\text{ m}$ thì mặt hồ là hình vuông và diện tích mặt hồ giảm được $20\text{ m}^2$. Tính chiều dài, chiều rộng của mảnh vườn hình chữ nhật.
\loigiai{
Gọi chiều rộng của mặt hồ ban đầu là $x\text{ (m)}$ và chiều dài của mặt hồ ban đầu là $y\text{ (m)}$ ($y > 5, x > 0, y > x$).\\
Khi giảm chiều dài $5\text{ m}$ và tăng chiều rộng $2\text{ m}$ thì mặt hồ trở thành hình vuông, ta có phương trình:
$y - 5 = x + 2 \Leftrightarrow y - x = 7 \Leftrightarrow y = x + 7$.\\
Diện tích mặt hồ ban đầu là $x \cdot y\text{ (m}^2\text{)}$.\\
Diện tích mặt hồ sau khi thay đổi kích thước là $(y - 5)(x + 2) = (x + 2)^2\text{ (m}^2\text{)}$.\\
Vì diện tích mặt hồ mới giảm $20\text{ m}^2$ so với ban đầu nên ta có phương trình:
$\begin{aligned}
xy - (x + 2)^2 = 20 &\Leftrightarrow x(x + 7) - (x^2 + 4x + 4) = 20 \\
&\Leftrightarrow x^2 + 7x - x^2 - 4x - 4 = 20 \\
&\Leftrightarrow 3x = 24 \Leftrightarrow x = 8\text{ (thỏa mãn)}.
\end{aligned}$\\
Suy ra chiều rộng của mặt hồ là $8\text{ m}$ và chiều dài của mặt hồ là $y = 8 + 7 = 15\text{ m}$.\\
Diện tích của mặt hồ là:
$S_{\text{hồ}} = 8 \cdot 15 = 120\text{ (m}^2\text{)}$.\\
Vì diện tích mặt hồ chiếm $3\%$ diện tích mảnh vườn nên diện tích mảnh vườn là:
$S_{\text{vườn}} = \dfrac{120}{3\%} = \dfrac{120}{0{,}03} = 4\,000\text{ (m}^2\text{)}$.\\
Gọi chiều rộng của mảnh vườn là $W\text{ (m)}$ ($W > 0$). Chiều dài của mảnh vườn là $2{,}5W\text{ (m)}$.\\
Ta có phương trình diện tích mảnh vườn:
$W \cdot 2{,}5W = 4\,000 \Leftrightarrow 2{,}5W^2 = 4\,000 \Leftrightarrow W^2 = 1\,600 \Rightarrow W = 40\text{ (m)}$.\\
Chiều dài của mảnh vườn là:
$2{,}5 \cdot 40 = 100\text{ (m)}$.\\
Vậy mảnh vườn có chiều rộng là $40\text{ m}$ và chiều dài là $100\text{ m}$.
}
\end{bt}

\begin{bt}%[9H9C29-3]
% nguon: Đề tham khảo Tuyển sinh 10 TP.HCM 2025-2026 – Quận 5 – Đề 2 – Bài 7
Cho đường tròn $(O, R)$. Từ điểm $A$ nằm ngoài đường tròn $(O)$, vẽ hai tiếp tuyến $AB$ và $AC$ đến đường tròn $(O)$ ($B$, $C$ là hai tiếp điểm). Vẽ đường kính $CD$ của đường tròn $(O)$.
\begin{enumerate}[a)]
    \item Chứng minh tứ giác $OBAC$ nội tiếp và $OA \parallel BD$.
    \item Gọi $H$ là giao điểm của $OA$ và $BC$. Chứng minh $OA \cdot OH = OB^2$ và $\widehat{ODH} = \widehat{ADB}$.
    \item Vẽ $BN$ song song với $AC$ ($N \in CD$). Trong trường hợp $OA = 2R$ tính diện tích của tứ giác $DBAC$ và độ dài $CN$ theo $R$.
\end{enumerate}
\loigiai{
\begin{enumerate}[a)]
    \item Vì $AB, AC$ là hai tiếp tuyến của $(O)$ nên $OB \perp AB$ và $OC \perp AC$, suy ra $\widehat{ABO} = \widehat{ACO} = 90^\circ$.\\
    Tứ giác $OBAC$ có tổng hai góc đối: $\widehat{ABO} + \widehat{ACO} = 90^\circ + 90^\circ = 180^\circ$.\\
    Vậy tứ giác $OBAC$ nội tiếp đường tròn đường kính $OA$.\\
    Mặt khác, $CD$ là đường kính của $(O)$ nên $\widehat{CBD} = 90^\circ$ (góc nội tiếp chắn nửa đường tròn), suy ra $BD \perp BC$.\\
    Theo tính chất hai tiếp tuyến cắt nhau, ta có $AB = AC$ và $OB = OC = R$, suy ra $OA$ là đường trung trực của đoạn thẳng $BC$, do đó $OA \perp BC$.\\
    Vì $BD \perp BC$ và $OA \perp BC$ nên $OA \parallel BD$.
    \item Xét tam giác $OBA$ vuông tại $B$ có đường cao $BH$ ($H$ là giao điểm của $OA$ và $BC$):\\
    Áp dụng hệ thức lượng trong tam giác vuông:
    $OA \cdot OH = OB^2 = R^2$.\\
    Vì $OD = R$ nên $OA \cdot OH = OD^2 \Leftrightarrow \dfrac{OH}{OD} = \dfrac{OD}{OA}$.\\
    Xét $\triangle ODH$ và $\triangle OAD$ có:
    $\left\{\begin{aligned} &\widehat{AOD}\text{ là góc chung} \\ &\dfrac{OH}{OD} = \dfrac{OD}{OA}\text{ (chứng minh trên)}. \end{aligned}\right.$\\
    Do đó $\triangle ODH \sim \triangle OAD$ (c-g-c), suy ra $\widehat{ODH} = \widehat{OAD}$.\\
    Mặt khác, do $OA \parallel BD$ nên hai góc so le trong bằng nhau: $\widehat{ADB} = \widehat{OAD}$.\\
    Từ đó suy ra $\widehat{ODH} = \widehat{ADB}$.
    \item Khi $OA = 2R$:\\
    Xét tam giác vuông $OBA$ vuông tại $B$ có cạnh huyền $OA = 2R$, cạnh góc vuông $OB = R$:\\
    Áp dụng định lí Pythagore:
    $AB = \sqrt{OA^2 - OB^2} = \sqrt{4R^2 - R^2} = R\sqrt{3}$.\\
    Ta có $\sin \widehat{BOA} = \dfrac{AB}{OA} = \dfrac{R\sqrt{3}}{2R} = \dfrac{\sqrt{3}}{2} \Rightarrow \widehat{BOA} = 60^\circ$.\\
    Hệ thức lượng trong tam giác vuông $OBA$:
    $\begin{aligned}
    &BH = \dfrac{OB \cdot AB}{OA} = \dfrac{R \cdot R\sqrt{3}}{2R} = \dfrac{R\sqrt{3}}{2} \Rightarrow BC = 2BH = R\sqrt{3}, \\
    &OH = \dfrac{OB^2}{OA} = \dfrac{R^2}{2R} = \dfrac{R}{2} \Rightarrow AH = OA - OH = 2R - \dfrac{R}{2} = \dfrac{3R}{2}.
    \end{aligned}$\\
    Tam giác $BCD$ vuông tại $B$ có $CD = 2R$ và $BC = R\sqrt{3}$, suy ra:
    $BD = \sqrt{CD^2 - BC^2} = \sqrt{4R^2 - 3R^2} = R$.\\
    Diện tích tam giác vuông $DBC$ là:
    $S_{\triangle DBC} = \dfrac{1}{2}BD \cdot BC = \dfrac{1}{2} \cdot R \cdot R\sqrt{3} = \dfrac{R^2\sqrt{3}}{2}$.\\
    Diện tích tam giác $ABC$ có đường cao $AH$ ứng với cạnh đáy $BC$ là:
    $S_{\triangle ABC} = \dfrac{1}{2}AH \cdot BC = \dfrac{1}{2} \cdot \dfrac{3R}{2} \cdot R\sqrt{3} = \dfrac{3R^2\sqrt{3}}{4}$.\\
    Vậy diện tích tứ giác $DBAC$ là:
    $S_{DBAC} = S_{\triangle DBC} + S_{\triangle ABC} = \dfrac{R^2\sqrt{3}}{2} + \dfrac{3R^2\sqrt{3}}{4} = \dfrac{5R^2\sqrt{3}}{4}$.\\
    Tính độ dài $CN$:\\
    Vì $AC$ là tiếp tuyến tại $C$ của $(O)$ nên $AC \perp CD$.\\
    Đề cho $BN \parallel AC$, suy ra $BN \perp CD$ tại điểm $N$.\\
    Do đó $BN$ là đường cao của tam giác $DBC$ vuông tại $B$.\\
    Áp dụng hệ thức lượng trong tam giác vuông $DBC$:
    $BC^2 = CD \cdot CN \Leftrightarrow (R\sqrt{3})^2 = 2R \cdot CN \Leftrightarrow 3R^2 = 2R \cdot CN \Rightarrow CN = \dfrac{3R}{2}$.\\
    Vậy diện tích tứ giác $DBAC$ là $\dfrac{5R^2\sqrt{3}}{4}$ và $CN = \dfrac{3R}{2}$.
\end{enumerate}
}
\end{bt}
